\scanpage{86}

% L1–2, top right
\margin{\uline{NB} Dans f), les conditions $\alpha)$ $\beta)$ $\gamma)$ ne font intervenir que $(\unsure{\tau_i|\mathrm T_i})$, non le choix de $s_0$, qui est précisé dans $\delta)\,\unsure{\eta|\gamma})\,\unsure{\varphi|\xi|\rho})$.}

% L2–7
e) Posant

% L3–7
\[
\begin{aligned}
s_0&=\text{unique pt de }X_0\\
s_1&=\tau_0s_0\\
s_2&=\tau_1s_1\\
&\ldots\\
s_{n+1}&=\tau_n s_n,
\end{aligned}
\]

% L4–8, upper right
\margin{\uline{NB}\quad (5) $s_i=\deleted{\unsure{\tau_i}}\tau_{i-1}\ldots\tau_0s_0$ $(1\leq i\leq n+1)$\\
\textcircled{6} donc si $j>i$ $\tau_js_i=(\tau_{i-1}\ldots\tau_0)\tau_js_0=s_i$ car $\tau_js_0=s_0$.}

% L3–10, left diagonal
% slanted
\margin{il y a une condition duale e$'$), où $\deleted{X_n}S_{n+1}=X_n$, $S_{\unsure{n|n-1}}=\tau_n X_n$, $S_{n-1}=\tau_{n-\unsure{2|1}}S_{n-1}$ \ldots\ $S_{\deleted{\unsure{0}}1}=\tau_0 X_0$, en exigeant $\bigwedge_{0\leq i\leq n+1}S_i=\varnothing$}

% L8–9
les $s_i$ sont \deleted{\unsure{affin.t}} projectiv$^{\mathrm{nt}}$ indépendants, i.e. engendrent l'espace $X$.

% L10
[ f) \deleted{Soit} \deleted{\unsure{$h$}} $h_i$ ($H_i$) le centre (l'antre) de $\tau_i$. On a :

% L10–17, right
\margin{$\alpha)$ $\bigvee h_i$ hyperplan $C$\\
$\beta)$ $\bigwedge H_i$ pt $c$ ($H_i$\unsure{/|,|)} indép.)\\
$\gamma)$ $h_i\notin H_{i+1}$ $(0\leq i\leq n-1)$\\
$\delta)$ $s_0\in\underbrace{H_1\cap\cdots\cap H_n}_{(\text{droite }D)}$\\
(NB il est plus fort a priori que $\tau_i s_0=s_0\quad1\leq i\leq n$)\\
$\unsure{\eta|\gamma})$ $s_0\notin H_0$ ou encore $s_0\ne c$ ($\{c\}=\bigwedge H_i$) ou encore $\tau_0s_0\ne s_0$\\
$\unsure{\varphi|\xi|\rho})$ $s_0\notin C$ i.e. $s_0\ne d$ où $\{c'\}=D\cap C$.}

% L11
\uline{Schéma de dém} On a les implications

% L12–14, implication scheme
\diagram{\begin{tikzcd}[ampersand replacement=\&,column sep=small,row sep=small]
a) \arrow[dr,Rightarrow,"{\mathrm{tr}}"] \& \& \& \& a)\\
\& b) \arrow[r,Leftrightarrow,"{\mathrm{tr}}"] \& c) \arrow[r,Rightarrow] \arrow[d,Leftrightarrow,"{\mathrm{tr}}"'] \& e) \arrow[ur,Rightarrow,"{\mathrm{tr}}"] \arrow[dr,Rightarrow,"{\mathrm{tr}}"'] \&\\
a') \arrow[ur,Rightarrow,"{\mathrm{tr}}"] \& \& d) \& \& a')
\end{tikzcd}}

% L15
la seule implication non triviale étant $d)\Longrightarrow e)$

% L16–18
On prouve l'indép. des $(s_0,\ldots,s_i)$ par \deleted{\unsure{réc}} sur $i$, en \deleted{Si elle est établie} prouvant en \unsure{même} temps

% L19
\uline{Cor} Sous les conditions précédentes, on a

% L20–21
\[
\left\{\begin{aligned}
X_i&=\operatorname{Lin}(s_0,\ldots,s_i)\\
&\deleted{\tau_j s_i=s_i\ \text{si }j\unsure{>|\ne}i}
\end{aligned}\right.\quad
\begin{aligned}&(0\leq i\leq n+1,\deleted{\unsure{S}}\\[-2pt]
 &\text{posant }X_{n+1}=X)\end{aligned}
\tag*{(7)}
\]

% L22–23
Si c'est prouvé pour \deleted{$i$, \unsure{savoir}, $\tau_i$} $i\leq i_0$, prouvons le pour $i=i_0+1$, i.e.

% L24–26
\[
\left\{\begin{aligned}
X_{i+1}&=\operatorname{Lin}(s_0,\ldots,s_i,s_{i+1})\\
&\quad\bigl(\operatorname{Lin}(X_i,s_{i+1})\\
\tau_j s_{i+1}&=s_{i+1}\quad\text{si }j>i+1
\end{aligned}\right.
\]

% L26–27, extension to right
\margin{$\longrightarrow$ \deleted{\unsure{vérifié par l'hyp. de récurrence|vérifié par (5)}} \unsure{est|car} $s_i\in X_i\subset X_{i+1}$}

% L27–31
Par déf $s_{i+1}=\tau_i s_i$, or $\tau_i X_{i+1}=X_{i+1}$ donc \uline{$s_{i+1}\in X_{i+1}$}. \deleted{\unsure{On|Reste}} \inserted{Il faut prouver qu'on} ne peut avoir $s_{i+1}\in X_i$ i.e. $\tau_i s_i\in X_i$, car \deleted{si} $X_i=\operatorname{lin}(X_{i-1},s_i)$ donc $\tau_iX_i=\operatorname{lin}(\underbrace{\tau_i X_{i-1}}_{X_{i-1}},\underbrace{\tau_i s_i}_{\unsure{\text{OK}}})$ qui serait égal à $X_i$, contraire à l'hyp. \unsure{OK|c)}

% archive sheet number: 85
